% Report on the third experiment, 23 November 2026.
%
%   .venv/bin/python v3/writeup/figures.py                       (figures and numbers)
%   cd v3/writeup && ~/.local/share/tectonic-0.17.0/tectonic report.tex
%
% Every number in the text is a macro from numbers.tex, which figures.py writes from the run
% files. Nothing is typed by hand, so text and figures cannot disagree.
%
% The rate check (the best repeat of every setting again at a third of the learning rate) is
% part of the report: finding 8, with a line on page one.

\documentclass[10pt,letterpaper]{article}
\usepackage[letterpaper,textwidth=6.4in,top=0.68in,bottom=0.86in,footskip=0.4in]{geometry}

% One family of typefaces for text and figures: IBM Plex. Serif for the text, Sans for headings,
% captions, tables and everything inside the figures, Mono for words of the made-up language.
\usepackage{fontspec}
\setmainfont{IBMPlexSerif}[Path=fonts/,Extension=.otf,UprightFont=*-Regular,ItalicFont=*-Italic,
  BoldFont=*-SemiBold,BoldItalicFont=*-BoldItalic]
\setsansfont{IBMPlexSans}[Path=fonts/,Extension=.otf,UprightFont=*-Regular,ItalicFont=*-Italic,
  BoldFont=*-SemiBold]
\setmonofont{IBMPlexMono}[Path=fonts/,Extension=.otf,UprightFont=*-Regular,BoldFont=*-SemiBold,Scale=0.9]
\linespread{1.09}

\usepackage{microtype}
\usepackage{graphicx}
\usepackage{float}
\usepackage{flafter}   % a figure never appears before the place where the text brings it up
\usepackage{needspace}
\usepackage{xcolor}
\usepackage{booktabs,array}
\usepackage{enumitem}
\usepackage{titlesec}
\usepackage[most]{tcolorbox}
\usepackage{fancyhdr}
\usepackage[font={sf,footnotesize},labelfont=bf,labelsep=period,skip=5pt,justification=raggedright,
  singlelinecheck=false]{caption}
\usepackage[hidelinks,pdftitle={Do more table columns let a model know more? Third experiment}]{hyperref}

% The colours of the figures (see the top of figures.py).
\definecolor{ink}{HTML}{1A1A19}
\definecolor{second}{HTML}{52514E}
\definecolor{muted}{HTML}{898781}
\definecolor{hair}{HTML}{C3C2B7}
\definecolor{pale}{HTML}{F4F3EF}
\color{ink}

\setlength{\parindent}{0pt}
\setlength{\parskip}{4.6pt plus 1pt minus 0.5pt}
\setlist{itemsep=2pt,topsep=1.5pt,parsep=0pt,partopsep=0pt,leftmargin=1.25em}
\raggedbottom
\widowpenalty=10000 \clubpenalty=10000
\emergencystretch=1.5em

% Headings. A finding is a full sentence under a small label; so is the part on the predictions,
% which stands beside the findings.
\titleformat{\section}{\sffamily\bfseries\fontsize{14.5}{18}\selectfont}{\textcolor{muted}{\thesection}}{0.55em}{}
\titlespacing*{\section}{0pt}{17pt plus 3pt minus 2pt}{3pt}
\titleformat{\subsection}{\sffamily\bfseries\fontsize{10.5}{13}\selectfont}{}{0pt}{}
\titlespacing*{\subsection}{0pt}{9pt plus 2pt minus 1pt}{0pt}
\newcommand{\smallcaps}[1]{{\sffamily\fontsize{7.6}{9}\selectfont\color{second}\addfontfeature{LetterSpace=7}#1}}
\newcommand{\labelled}[2]{%
  \vspace{12pt plus 3pt minus 2pt}%
  \smallcaps{#1}\par\nopagebreak\vspace{-2pt}%
  {\sffamily\bfseries\fontsize{12.5}{15.5}\selectfont\raggedright #2\par}%
  \nopagebreak\vspace{1pt}}
\newcounter{finding}
\newcommand{\finding}[1]{\refstepcounter{finding}\labelled{FINDING \thefinding}{#1}}

\newcommand{\s}[1]{\texttt{#1}}                  % a word or sentence of the made-up language, or a file
\newcommand{\term}[1]{\textbf{\textit{#1}}}      % a term at the place where it is explained
\newcommand{\lead}[1]{\textbf{#1}}               % the opening words of a point
\newcommand{\point}[1]{{\color{ink}\bfseries #1}} % the first sentence of a caption: the point
\newcommand{\see}[1]{\mbox{\sffamily\footnotesize\color{second}(#1)}} % where an answer on page one is shown
\newcommand{\everyrun}{pages \pageref{page:everyfirst} to~\pageref{page:everylast}} % where every run is shown
\captionsetup{textfont={color=second}}
\setlength{\intextsep}{9pt plus 2pt minus 1pt}
\setlength{\textfloatsep}{11pt plus 2pt minus 1pt}
\renewcommand{\topfraction}{0.9}\renewcommand{\bottomfraction}{0.9}\renewcommand{\textfraction}{0.05}
\renewcommand{\floatpagefraction}{0.85}\setcounter{topnumber}{2}\setcounter{totalnumber}{4}
% A figure at full width with its caption underneath, exactly where it is called: the report is
% laid out so that each figure stands with the text that reads it.
\newcommand{\fighere}[3]{\begin{figure}[H]\includegraphics[width=\linewidth]{figures/#1.pdf}%
  \caption{#3}\label{#2}\end{figure}}

\pagestyle{fancy}\fancyhf{}\renewcommand{\headrulewidth}{0pt}
\fancyfoot[L]{\sffamily\fontsize{7.6}{9}\selectfont\color{muted}Do more table columns let a model know more?
  \enspace·\enspace Third experiment \enspace·\enspace 23 November 2026}
\fancyfoot[R]{\sffamily\fontsize{7.6}{9}\selectfont\color{second}\thepage}

% The four pages that show every run lie on their side, so that a page holds all 720 runs of one
% number of table columns. \pageshape starts a new page of the given paper width and height
% with a text block of the given width and height, moved up by the last length.
\makeatletter
\newcommand{\pageshape}[5]{%
  \clearpage
  \pdfpagewidth=#1 \pdfpageheight=#2 \paperwidth=#1 \paperheight=#2
  \textwidth=#3 \hsize=#3 \linewidth=#3 \columnwidth=#3 \headwidth=#3
  \oddsidemargin=\dimexpr(#1-#3)/2-1in\relax \evensidemargin=\oddsidemargin
  \addtolength{\topmargin}{-#5}%
  \textheight=#4 \vsize=#4 \global\@colht=#4 \global\@colroom=#4}
\makeatother

\newcommand{\ComputeUnits}{282}
\newcommand{\StageStart}{08:56}
\newcommand{\StageEnd}{14:35}
\newcommand{\SharingStart}{12:48}
\newcommand{\MovedLateGroups}{twelve}
\newcommand{\MovedLateMachines}{five}
\newcommand{\MachinesAtOnce}{ten}
\newcommand{\MachinesPerType}{three}
\newcommand{\PlannedSteps}{357}
\newcommand{\SentencesPerStep}{512}
\newcommand{\Layers}{2}
\newcommand{\Heads}{2}
\newcommand{\RunsAll}{2,880}
\newcommand{\Settings}{960}
\newcommand{\RunsPerUsePoint}{18}
\newcommand{\SettingsPerColumns}{240}
\newcommand{\FirstLook}{40,000}
\newcommand{\Limit}{640,000}
\newcommand{\LearnedBar}{99}
\newcommand{\RunsPerGroup}{24}
\newcommand{\Groups}{120}
\newcommand{\Rate}{0.003}
\newcommand{\StepsMillions}{708}
\newcommand{\AttentionTimes}{58}
\newcommand{\OutputManySmall}{2,970}
\newcommand{\OutputManyLarge}{21,450}
\newcommand{\SizeGridLow}{9}
\newcommand{\SizeGridHigh}{26}
\newcommand{\BodyGridLow}{1.2}
\newcommand{\BodyGridHigh}{3.7}
\newcommand{\BodyFewLow}{1.4}
\newcommand{\BodyFewHigh}{3.7}
\newcommand{\BodyManyLow}{1.2}
\newcommand{\BodyManyHigh}{1.4}
\newcommand{\HeldCount}{two}
\newcommand{\NotHeldCount}{three}
\newcommand{\COneComparisons}{720}
\newcommand{\COneFalls}{0}
\newcommand{\COneLargestFall}{0.1}
\newcommand{\COneWentUp}{292}
\newcommand{\COneStayedLevel}{428}
\newcommand{\COneCases}{240}
\newcommand{\COneGained}{124}
\newcommand{\COneNeeded}{60}
\newcommand{\COneNotGained}{116}
\newcommand{\COneLargestGain}{61}
\newcommand{\COneRoom}{124}
\newcommand{\RoomBar}{90}
\newcommand{\COneNoRoom}{101}
\newcommand{\COneByChance}{three}
\newcommand{\COneWithoutChance}{121}
\newcommand{\COneRunFall}{0.2}
\newcommand{\COneRunLow}{113}
\newcommand{\COneRunHigh}{133}
\newcommand{\TypicalEqual}{40, 56, 78 and 100}
\newcommand{\TypicalUnequal}{33, 44, 60 and 78}
\newcommand{\TypicalGapLow}{10}
\newcommand{\TypicalGapHigh}{22}
\newcommand{\PooledCOneGained}{61}
\newcommand{\PooledCOneCases}{120}
\newcommand{\PooledCTwoGained}{9}
\newcommand{\PooledCTwoCases}{80}
\newcommand{\CTwoCases}{160}
\newcommand{\CTwoGained}{19}
\newcommand{\CTwoAllowed}{15}
\newcommand{\CTwoPercent}{12}
\newcommand{\CTwoOver}{four}
\newcommand{\CTwoAtEight}{twelve}
\newcommand{\CTwoAtSixteen}{five}
\newcommand{\CTwoAtThirtyTwo}{two}
\newcommand{\CTwoAtSixtyFour}{zero}
\newcommand{\CTwoLargest}{27.7}
\newcommand{\CTwoLargestWhere}{8 columns, 128 bits, 64 creatures, equal use}
\newcommand{\CTwoLargestFrom}{55}
\newcommand{\CTwoLargestTo}{83}
\newcommand{\CTwoByChance}{one}
\newcommand{\CTwoWithoutChance}{18}
\newcommand{\CTwoRoom}{39}
\newcommand{\CTwoNoRoom}{121}
\newcommand{\CTwoRoomFew}{9}
\newcommand{\CTwoRoomFewGained}{9}
\newcommand{\CTwoRoomMid}{12}
\newcommand{\CTwoRoomMidGained}{7}
\newcommand{\CTwoRoomMany}{18}
\newcommand{\CTwoRoomManyGained}{3}
\newcommand{\CTwoRunLow}{10}
\newcommand{\CTwoRunHigh}{26}
\newcommand{\DoubledCases}{120}
\newcommand{\DoubledMore}{40}
\newcommand{\DoubledTied}{80}
\newcommand{\BodyTimes}{32}
\newcommand{\SameSizePairs}{19}
\newcommand{\SameSizeLanguages}{nine}
\newcommand{\SameSizePercent}{15}
\newcommand{\SameSizeHigher}{19}
\newcommand{\SameSizeByTwo}{10}
\newcommand{\SameSizeByFive}{four}
\newcommand{\SameSizeLargest}{23}
\newcommand{\SameSizeSmaller}{five}
\newcommand{\SameSizeAllRuns}{18}
\newcommand{\SameSizeMostCreatures}{256}
\newcommand{\SizeFirstWide}{8,538}
\newcommand{\SizeFirstNarrow}{9,490}
\newcommand{\SizeFirstWideScore}{99}
\newcommand{\SizeFirstNarrowScore}{100}
\newcommand{\SizeFirstGap}{1.4}
\newcommand{\SizeFirstPercent}{11}
\newcommand{\SizeSecondWide}{8,538}
\newcommand{\SizeSecondNarrow}{9,490}
\newcommand{\SizeSecondWideScore}{69}
\newcommand{\SizeSecondNarrowScore}{92}
\newcommand{\SizeSecondGap}{23.2}
\newcommand{\SizeSecondPercent}{11}
\newcommand{\SizeThirdWide}{18,938}
\newcommand{\SizeThirdNarrow}{30,130}
\newcommand{\SizeThirdWideScore}{42}
\newcommand{\SizeThirdNarrowScore}{55}
\newcommand{\SizeThirdGap}{12.3}
\newcommand{\SizeThirdPercent}{59}
\newcommand{\SizeManyClosest}{1.4}
\newcommand{\AOneCases}{480}
\newcommand{\AOneLearned}{432}
\newcommand{\AOneMissed}{48}
\newcommand{\AOneAllThree}{400}
\newcommand{\AOneBelow}{288}
\newcommand{\AOneAtOne}{192}
\newcommand{\AOneFactsKnown}{475}
\newcommand{\AOneQuestionsShort}{43}
\newcommand{\AOneFactsShort}{five}
\newcommand{\AOneQuestionsOnly}{32}
\newcommand{\AOneRuns}{1,440}
\newcommand{\AOneRunsLearned}{1,292}
\newcommand{\AOneShortByKind}{three}
\newcommand{\AOneShortOthers}{two}
\newcommand{\AOneShortOthersBy}{3.4 and 0.03}
\newcommand{\AOneMissedEight}{30}
\newcommand{\ATwoCases}{72}
\newcommand{\ATwoReached}{0}
\newcommand{\ATwoHighest}{86.5}
\newcommand{\ATwoBar}{95}
\newcommand{\ATwoHighestRun}{87.3}
\newcommand{\ATwoRuns}{216}
\newcommand{\ATwoRunsLimit}{166}
\newcommand{\ATwoTopGainLow}{2.1}
\newcommand{\ATwoTopGainHigh}{3.6}
\newcommand{\AThreeCases}{960}
\newcommand{\AThreeWithin}{891}
\newcommand{\AThreeApart}{69}
\newcommand{\AThreeMoreHigher}{64}
\newcommand{\AThreeFewerHigher}{five}
\newcommand{\AThreeByChance}{nine}
\newcommand{\AThreeWithoutChance}{60}
\newcommand{\StuckTile}{93}
\newcommand{\SameLoadEqual}{90, 94, 98 and 100}
\newcommand{\SameLoadMidEqualLow}{98}
\newcommand{\SameLoadMidEqualHigh}{100}
\newcommand{\SameLoadUnequal}{88, 86, 84 and 78}
\newcommand{\SameLoadMidUnequalLow}{90}
\newcommand{\SameLoadMidUnequalHigh}{97}
\newcommand{\SameLoadRuns}{144}
\newcommand{\SameLoadRunsLimit}{134}
\newcommand{\CreaturesExample}{100, 98, 78 and 67}
\newcommand{\FewCreaturesLowest}{80}
\newcommand{\PlateauRuns}{12}
\newcommand{\PlateauRunsWord}{twelve}
\newcommand{\PlateauHere}{eight}
\newcommand{\PlateauGroup}{18}
\newcommand{\PlateauScore}{81.25}
\newcommand{\PlateauScoreEight}{90.6}
\newcommand{\PlateauFewer}{two}
\newcommand{\PlateauMoreBits}{two}
\newcommand{\PlateauExample}{81.2, 99.9 and 81.2}
\newcommand{\PlateauBeside}{three}
\newcommand{\PlateauBodies}{six}
\newcommand{\PlateauCarried}{three}
\newcommand{\PlateauCarriedLearned}{one}
\newcommand{\PlateauCarriedNot}{two}
\newcommand{\PlateauOthers}{eleven}
\newcommand{\PlateauWiderRuns}{2,160}
\newcommand{\SpreadMiddle}{0.02}
\newcommand{\SpreadOverFive}{52}
\newcommand{\NoGainRuns}{574}
\newcommand{\NoGainKnew}{424}
\newcommand{\NoGainBelow}{150}
\newcommand{\NoGainMiddle}{0.5}
\newcommand{\NoGainOverHalf}{79}
\newcommand{\NoGainEight}{130}
\newcommand{\NoGainLossFall}{0.008}
\newcommand{\NoGainLoss}{1.8}
\newcommand{\Marked}{288}
\newcommand{\MarkedNoGain}{82}
\newcommand{\MarkedAtLastLook}{ten}
\newcommand{\Carried}{72}
\newcommand{\CarriedKnew}{55}
\newcommand{\CarriedRoom}{17}
\newcommand{\CarriedLevel}{eleven}
\newcommand{\CarriedGained}{six}
\newcommand{\CarriedGainedDigits}{6}
\newcommand{\CarriedPercent}{8}
\newcommand{\CarriedOthers}{five}
\newcommand{\CarriedOthersLow}{1.0}
\newcommand{\CarriedOthersHigh}{2.3}
\newcommand{\CarriedOneDoubling}{15}
\newcommand{\JumpFrom}{81.2}
\newcommand{\JumpTo}{100}
\newcommand{\JumpGain}{18.8}
\newcommand{\JumpQuestions}{99}
\newcommand{\JumpStepsAfter}{32,000}
\newcommand{\CarriedQuestionsGained}{58}
\newcommand{\CarriedLight}{22}
\newcommand{\CarriedLightLearned}{11}
\newcommand{\LimitRuns}{757}
\newcommand{\LimitStillGaining}{582}
\newcommand{\LimitGainLow}{1.4}
\newcommand{\LimitGainHigh}{1.8}
\newcommand{\AtLineFactsLow}{100}
\newcommand{\AtLineFactsHigh}{100}
\newcommand{\AtLineQuestionsLow}{98}
\newcommand{\AtLineQuestionsHigh}{100}
\newcommand{\AtLineNewLow}{93}
\newcommand{\AtLineNewHigh}{97}
\newcommand{\AtLineNotLow}{81}
\newcommand{\AtLineNotHigh}{95}
\newcommand{\HeavyFactsEqual}{100}
\newcommand{\HeavyQuestionsEqual}{77}
\newcommand{\HeavyNewEqual}{72}
\newcommand{\HeavyNotEqual}{46}
\newcommand{\HeavyFactsUnequal}{78}
\newcommand{\HeavyQuestionsUnequal}{62}
\newcommand{\HeavyNewUnequal}{67}
\newcommand{\HeavyNotUnequal}{47}
\newcommand{\EarlierSettings}{120}
\newcommand{\EarlierApart}{2.9}
\newcommand{\EarlierHeavy}{18}
\newcommand{\EarlierHeavyLow}{0}
\newcommand{\EarlierHeavyHigh}{6.2}
\newcommand{\PracticeRuns}{240}
\newcommand{\PracticeSettings}{80}
\newcommand{\PracticeRates}{0.003, 0.01 and 0.03}
\newcommand{\PracticeSteps}{40,000}
\newcommand{\PracticeEightLow}{78.5}
\newcommand{\PracticeSixteenLow}{88.4}
\newcommand{\PracticeThirtyTwoLow}{95.8}
\newcommand{\PracticeSixtyFourLow}{98.3}
\newcommand{\PracticeEightMid}{79.4}
\newcommand{\PracticeSixteenMid}{89.5}
\newcommand{\PracticeThirtyTwoMid}{95.5}
\newcommand{\PracticeSixtyFourMid}{97.7}
\newcommand{\PracticeEightHigh}{77.9}
\newcommand{\PracticeSixteenHigh}{83.0}
\newcommand{\PracticeThirtyTwoHigh}{86.6}
\newcommand{\PracticeSixtyFourHigh}{80.0}
\newcommand{\PracticeRows}{0.003 & 90.2 & 75.3 & 78.5 & 88.4 & 95.8 & 98.3 \\ 0.01 & 90.5 & 73.0 & 79.4 & 89.5 & 95.5 & 97.7 \\ 0.03 & 81.9 & 56.0 & 77.9 & 83.0 & 86.6 & 80.0 \\}
\newcommand{\Machines}{12}
\newcommand{\MachineTypes}{five}
\newcommand{\MachineList}{two T4, three L4, three A100, three G4 and one H100}
\newcommand{\MachinesPaused}{three}
\newcommand{\LauncherRestarts}{ten}
\newcommand{\GroupFailures}{one}
\newcommand{\GroupsOnSeveralMachines}{57}
\newcommand{\GroupsOnSeveralTypes}{49}
\newcommand{\NextBodyLow}{320}
\newcommand{\NextBodyHigh}{520}
\newcommand{\RcRate}{0.001}
\newcommand{\RcRuns}{960}
\newcommand{\RcTies}{544}
\newcommand{\RcHigher}{13}
\newcommand{\RcLower}{19}
\newcommand{\RcWithin}{928}
\newcommand{\RcUpEight}{29}
\newcommand{\RcDownEight}{86}
\newcommand{\RcUpSixteen}{12}
\newcommand{\RcDownSixteen}{39}
\newcommand{\RcUpThirtyTwo}{16}
\newcommand{\RcDownThirtyTwo}{10}
\newcommand{\RcUpSixtyFour}{9}
\newcommand{\RcDownSixtyFour}{0}
\newcommand{\RcDownSixtyFourWord}{none}
\newcommand{\RcLostMiddle}{2.7}
\newcommand{\RcTypicalEqual}{39, 54, 78 and 99}
\newcommand{\RcTypicalUnequal}{34, 44, 61 and 82}
\newcommand{\RcQuestionsHigher}{52}
\newcommand{\RcQuestionsLower}{172}
\newcommand{\RcLimitChosen}{249}
\newcommand{\RcLimitThird}{305}
\newcommand{\RcLimitGainChosen}{1.6}
\newcommand{\RcLimitGainThird}{2.0}
\newcommand{\RcCTwoChosen}{17}
\newcommand{\RcCTwoThird}{13}
\newcommand{\RcTopSixtyFour}{226}
\newcommand{\RcPerColumns}{240}
\newcommand{\CTwoSingleLow}{15}
\newcommand{\CTwoSingleHigh}{20}
\newcommand{\CTwoSameRepeatLow}{17}
\newcommand{\CTwoSameRepeatHigh}{20}
\newcommand{\RcCTwoByChance}{one}
\newcommand{\RcStart}{14:35}
\newcommand{\RcEnd}{16:47}
\newcommand{\RcMachines}{seven}
\newcommand{\RcMachineList}{three A100, three G4 and one H100}
\newcommand{\RcGroupFailures}{fifteen}
\newcommand{\RcTenth}{It was not run: the machines were given back to save compute units.}
\newcommand{\MovedLateGroupsCap}{Twelve}
\newcommand{\MovedLateMachinesCap}{Five}
\newcommand{\MachinesAtOnceCap}{Ten}
\newcommand{\MachinesPerTypeCap}{Three}
\newcommand{\HeldCountCap}{Two}
\newcommand{\NotHeldCountCap}{Three}
\newcommand{\COneByChanceCap}{Three}
\newcommand{\CTwoOverCap}{Four}
\newcommand{\CTwoAtEightCap}{Twelve}
\newcommand{\CTwoAtSixteenCap}{Five}
\newcommand{\CTwoAtThirtyTwoCap}{Two}
\newcommand{\CTwoAtSixtyFourCap}{Zero}
\newcommand{\CTwoByChanceCap}{One}
\newcommand{\SameSizeLanguagesCap}{Nine}
\newcommand{\SameSizeByFiveCap}{Four}
\newcommand{\SameSizeSmallerCap}{Five}
\newcommand{\AOneFactsShortCap}{Five}
\newcommand{\AOneShortByKindCap}{Three}
\newcommand{\AOneShortOthersCap}{Two}
\newcommand{\AThreeFewerHigherCap}{Five}
\newcommand{\AThreeByChanceCap}{Nine}
\newcommand{\PlateauRunsWordCap}{Twelve}
\newcommand{\PlateauHereCap}{Eight}
\newcommand{\PlateauFewerCap}{Two}
\newcommand{\PlateauMoreBitsCap}{Two}
\newcommand{\PlateauBesideCap}{Three}
\newcommand{\PlateauBodiesCap}{Six}
\newcommand{\PlateauCarriedCap}{Three}
\newcommand{\PlateauCarriedLearnedCap}{One}
\newcommand{\PlateauCarriedNotCap}{Two}
\newcommand{\PlateauOthersCap}{Eleven}
\newcommand{\MarkedAtLastLookCap}{Ten}
\newcommand{\CarriedLevelCap}{Eleven}
\newcommand{\CarriedGainedCap}{Six}
\newcommand{\CarriedOthersCap}{Five}
\newcommand{\MachineTypesCap}{Five}
\newcommand{\MachinesPausedCap}{Three}
\newcommand{\LauncherRestartsCap}{Ten}
\newcommand{\GroupFailuresCap}{One}
\newcommand{\RcDownSixtyFourWordCap}{None}
\newcommand{\RcCTwoByChanceCap}{One}
\newcommand{\RcMachinesCap}{Seven}
\newcommand{\RcGroupFailuresCap}{Fifteen}


\begin{document}

% ------------------------------------------------------------------------------------------------
% Page one: the question, the answer, the picture
% ------------------------------------------------------------------------------------------------

\smallcaps{THIRD EXPERIMENT \enspace·\enspace REPORT OF 23 NOVEMBER 2026}\par
\vspace{4pt}
{\sffamily\bfseries\fontsize{21}{25}\selectfont Do more table columns let a model know more?\par}
\vspace{4pt}
{\sffamily\fontsize{10.5}{14}\selectfont\color{second}\raggedright The claim under test is the owner's, of
22~November: ``the thing of interest here is not the width of the model but the dimensions in which the token
is represented. simply, there are more weights to help distinguish the tokens as the dimensions increase,
despite no change in width.''\par}
\vspace{5pt}{\color{hair}\hrule height 0.6pt}\vspace{6pt}

\lead{The test.} A model keeps a row of numbers for every word (every token) in its \term{word table}. The
length of that row, the number of \term{table columns}, is what the claim calls the dimensions of a token. The
experiment varies it (8 to 64) and four other dials, every combination three times, \RunsAll{} runs: the
\term{body width} (the working part of each layer, 4 to 128), the \term{bits} to know per creature (8 to
128), the creatures (16 to 1,024) and the \term{word use} (equal or unequal). The score is \term{facts known},
0 to 100.

\begin{tcolorbox}[enhanced,colback=pale,colframe=pale,boxrule=0pt,arc=0pt,left=11pt,right=11pt,top=6pt,
  bottom=7pt,borderline west={2.2pt}{0pt}{ink},before skip=7pt,after skip=2pt]
\smallcaps{THE ANSWER}\par
\vspace{0.5pt}
\begin{enumerate}[leftmargin=1.35em,itemsep=2.5pt,topsep=1pt,label=\textbf{\arabic*.}]
\item \lead{More columns, more facts known: this half of the claim held (prediction C1).} With all else
      fixed, no score fell by 5 points from one column count to the next (\COneFalls{} of \COneComparisons{}
      comparisons), and 64 columns scored 10 or more above 8 in \COneGained{} of \COneCases{} cases;
      \COneNeeded{} were asked for. \see{finding~\ref{f:columns}}
\item \lead{A wider body helped far less, yet more than the claim allowed: its other half narrowly did not hold
      (C2).} Body width 128 scored 10 points or more above body width 4 in \CTwoGained{} of \CTwoCases{}
      settings; the prediction allowed fewer than a tenth, at most \CTwoAllowed. \see{finding~\ref{f:body}}
\item \lead{Size alone does not explain it.} Where a model with 16 columns had about as many weights as one
      with 8, it knew more: in all \SameSizePairs{} such pairs. But more columns also enlarge parts of every
      layer, and which part does it, the runs cannot say. \see{finding~\ref{f:size}}
\item \lead{Of the agent's three predictions, one held and two did not.} They are about bits per creature
      divided by columns. A2 held: at 4 or more, none of \ATwoCases{} settings with 1,024 creatures reached
      \ATwoBar. A1 did not: at 1 or less all were to end as ``learned'' (facts and questions both at
      \LearnedBar); \AOneLearned{} of \AOneCases{} did, and in \AOneQuestionsShort{} of the other \AOneMissed{}
      the facts were at \LearnedBar, the questions short. A3 did not: equal bits per column were to score
      within 10 points; \AThreeApart{} of \AThreeCases{} pairs did not. \see{finding~\ref{f:line}}
\item \lead{A third of the learning rate changed little.} At that rate \RcWithin{} of \RcRuns{} settings
      ended within 5 points of their score, and four verdicts stayed. C2's count kept to its line of
      \CTwoAllowed: \RcCTwoChosen{} at the chosen rate and \RcCTwoThird{} at a third of it on single runs,
      \CTwoSingleLow{} to \CTwoSingleHigh{} with other repeats. \see{finding~\ref{f:rate}}
\end{enumerate}
\end{tcolorbox}

\fighere{fig01_answer}{fig:answer}{\point{The more table columns, the more facts a model knows; the lines part
as the creatures or the bits grow, and lie lower with unequal use.} A point: 6 body widths, 3 repeats each.
Every run: \everyrun.}

% ------------------------------------------------------------------------------------------------
\section{How the experiment is built}
% ------------------------------------------------------------------------------------------------

\subsection{The language and the test}

The language and the test are those of the second experiment, whose report shows both at length. Each made-up
creature has a fixed set of facts, and each fact is one of four possible answers (Figure~\ref{fig:language}).
Choosing one of four takes two yes-or-no choices, so a fact is worth 2 \term{bits}. A model is never shown the
facts. It only reads sentences about one to four creatures: statements (\s{foxes and crows are black~.}),
``not'' statements (\s{foxes do not eat seeds~.}) and questions with their answer (\s{are foxes and bears
black~?~no}). After training it sits a fixed test of four parts. Facts known is the main score, and the one that all
five predictions are about.

\fighere{fig02_language}{fig:language}{\point{A creature is a word with a set of facts, and a model is tested
on them in four ways.} Left: one creature in the language with 16 creatures, its true answers outlined. Right:
in each part of the test the model is given the words in the pale boxes and must supply the outlined one; for
``not'' known the true answer must get the lowest score of the four. A question score is the share of ``yes''
questions answered right plus the share of ``no'' questions answered right, minus 100, so always giving the
same answer scores 0.}

\subsection{The model, and what a table column is}

The model is a small \term{transformer}, the kind of model behind today's language systems. It reads a sentence
word by word and guesses the next word. A model is made of numbers, its \term{weights}. \term{Training} is this
guessing, repeated: after each round the weights are nudged so that the right words would have scored a little
higher. One such nudge is a \term{training step}; the model reads \SentencesPerStep{} sentences for each. This
report calls the weights the \term{trained numbers}.

The \term{word table} holds a row of numbers for every word, and these rows are trained with everything else. A
\term{table column} is one number of every row: a table with 8 columns gives each word 8 numbers. The
\term{body} is two layers that turn the rows of the words read so far into a guess. Each layer has an attention
part, which lets a word gather information from the words before it, and a middle part made of \term{neurons}:
small units that each add up what comes in and pass the result on. The \term{body width} is the number of
neurons in each middle part.

\lead{More columns change three things at once} (Figure~\ref{fig:columns}). A word's row is not only looked up.
It is also what every layer reads, and what every layer adds its work to. So the number of columns sets the size
of everything that touches a row:

\begin{enumerate}[leftmargin=1.5em,itemsep=1pt,topsep=1pt]
\item the row itself, for every word;
\item the attention parts: 8 times the columns give them about \AttentionTimes{} times the trained numbers;
\item the input to each middle part, and its output.
\end{enumerate}

The last part of the model follows the columns too: the \term{output part}, which turns the result into a score
for every word the model can answer with. Only the body width stays as it is. The plan says so, and the owner
chose it knowingly: ``if it is observed that this is the case, then we can narrow it further if that is
possible.'' So these runs can show \emph{that} more columns help. They cannot show \emph{which} of these parts
does it.

\fighere{fig03_columns}{fig:columns}{\point{More table columns make every word's row longer, the attention parts
larger and the input to the middle parts longer, all together.} The numbers count everything that training can
change. The output part keeps a row of weights for every word the model can answer with (at 128 bits per
creature, \OutputManySmall{} trained numbers with 8 columns and \OutputManyLarge{} with 64). The rest is a row
for each place in a sentence and a few numbers that keep values in range.}

\subsection{What was varied, and what was held the same}

Five dials were varied (Figure~\ref{fig:dials}). One combination of the five is a \term{setting}; there are
\Settings. Each setting was trained three times, its \term{repeats}: a repeat starts from other random numbers
in the body. Where the report gives one score for a setting, it is \term{the middle of its repeats}: the one of
the three with a higher and a lower beside it. On \everyrun{} all three are shown for every setting.

\fighere{fig04_dials}{fig:dials}{\point{Five dials, \Settings{} settings, \RunsAll{} runs, and everything else
the same.} The small bars beside the two word uses show how often the eight most used creatures appear. With
unequal use the second appears half as often as the first, the third a third as often, and so on, much as words
do in real text.}

Everything else was the same in every run, apart from the machines the runs trained on (Appendix~B). The owner
asked for this after the second experiment, whose stages had each been trained differently and could not be
compared cleanly.

\begin{itemize}
\item \lead{One learning rate} (how large a nudge each training step makes): \Rate, chosen by a fixed rule from
      \PracticeRuns{} practice runs before the \term{counted runs} began, the \RunsAll{} runs that are the
      results (Appendix~A). Finding~\ref{f:rate} tries a lower one.
\item \lead{One text.} Every model of a language reads exactly the same sentences in the same order.
\item \lead{One starting word table,} with 64 columns. A model with fewer columns starts from its first
      columns, so a word starts from the same numbers in every model.
\item \lead{\Layers{} layers and \Heads{} attention heads} at every size (an attention part is split into heads
      that look back separately).
\end{itemize}

\subsection{How a run ends}

Only the length of training differs between runs, and one rule sets it (Figure~\ref{fig:rule}). A fixed length
had left the hardest settings of the second experiment still rising at the end, so that ``cannot hold it'' could
not be told from ``has not learned it yet''. Here a run is looked at after \FirstLook{} steps and after each
doubling, up to \Limit. It ends at the first look where one of three things is true:

\begin{itemize}
\item \lead{learned:} facts known and questions about one creature are both \LearnedBar{} or more;
\item \lead{no more gain:} facts known rose by less than 1 point over the last doubling, and by less than 1
      point over the doubling before it. This also ends a run that knows the facts but is still short on the
      questions, because its facts cannot rise;
\item \lead{the limit:} it has reached \Limit{} steps.
\end{itemize}

A run's result is its scores at that point. The program applies the rule, not judgement. To measure how often
the rule stops a run too early, every tenth run was marked beforehand to be \term{carried on}: if it ended as
``no more gain'', it was trained for two more doublings anyway (finding~\ref{f:rule}).

\fighere{fig05_rule}{fig:rule}{\point{A run ends when it has learned, when it has stopped gaining, or at
\Limit{} steps; nearly half of all runs had learned at the first look.} Left: four runs with 1,024 creatures,
body width 16 and equal use, each drawn up to the look that ended it. Right: the \RunsAll{} runs by reason and
look. Three in four of the runs ended as ``no more gain'' knew the facts (\LearnedBar{} or more) and were short
only on the questions about one creature.}

% ------------------------------------------------------------------------------------------------
\clearpage  % the findings start on a new page, with the five predictions
\section{What the runs showed}
% ------------------------------------------------------------------------------------------------

\labelled{BEFORE THE RUNS}{Five predictions, fixed on 22 November}

The plan fixed five predictions before any run (Figure~\ref{fig:predictions}). C1 and C2 are the two halves of
the owner's claim. A1 to A3 are the agent's reading of the second experiment: what decides is bits per creature
divided by columns, the bits each number of a word's row must hold. In all five, ``score'' is facts known
where the rule ended a run, and a setting's score is the middle of its three repeats.

\fighere{fig06_predictions}{fig:predictions}{\point{\HeldCountCap{} of the five predictions held, C1 and A2; C2
missed narrowly, and A1 and A3 did not hold.} The dark bar is what the runs gave; the thin line under it spans
the counts the prediction allowed, with an upright mark at its limit. A setting counts as ``learned'' when at
least two of its three repeats ended so.}

\finding{More table columns, more facts known, wherever there was room}\label{f:columns}

Figure~\ref{fig:bycolumns}, at the top of the next page, is prediction C1 as a picture. From one column count
to the next larger one, no setting's score fell by more than \COneLargestFall{} points. Of the \COneComparisons{}
comparisons, \COneWentUp{} went up by more than half a point and \COneStayedLevel{} stayed within half a point,
all of them at \LearnedBar{} or more. From 8 to 64 columns the score rose by 10 points or more in \COneGained{}
of the \COneCases{} cases, by as much as \COneLargestGain{} points: exactly the cases that had room for such a
gain (\RoomBar{} or less with 8 columns). So C1 held, with about twice the cases it asked for;
\COneByChance{} of them are chance cases (finding~\ref{f:plateau}).

\lead{Run by run it is as firm:} no run ended more than \COneRunFall{} points below any run of the same case
with the next smaller number of columns, and with the worst repeat at 64 columns set against the best at 8,
\COneRunLow{} cases still rise by 10 points or more.

\clearpage
\fighere{fig07_by_columns}{fig:bycolumns}{\point{No line goes down: with everything else fixed, no model with
more table columns knew clearly fewer facts.} The \COneCases{} lines are the \COneCases{} cases of C1, their
\COneComparisons{} segments its \COneComparisons{} comparisons. The number in a panel counts its lines that end 10 points
or more higher at 64 columns than at 8. Lines that start at 81: finding~\ref{f:plateau}.}

\finding{A wider body helped far less, but more than the claim allowed}\label{f:body}

\fighere{fig08_body_gain}{fig:bodygain}{\point{Body width 128 scored 10 points or more above body width 4 in
\CTwoGained{} of \CTwoCases{} settings, \CTwoOver{} more than C2 allowed; nearly all of them have 8 or 16
columns and more bits than columns.} A tile is one setting of C2; its number is the gain from body width 4 to
128, each the middle of three repeats. The outlined tile left of the stepped line is a chance case
(finding~\ref{f:plateau}).}

So C2 did not hold: \CTwoGained{} of \CTwoCases{} settings (\CTwoPercent{} percent), where at most
\CTwoAllowed{} were allowed (Figure~\ref{fig:bodygain}). \CTwoAtEightCap{} of them have 8 columns,
\CTwoAtSixteen{} have 16, \CTwoAtThirtyTwo{} have 32 and none has 64. The largest gain was \CTwoLargest{}
points, from \CTwoLargestFrom{} to \CTwoLargestTo{} (\CTwoLargestWhere).

\begin{itemize}
\item \lead{Where there was room, the body often mattered.} Only \CTwoRoom{} of the \CTwoCases{} settings had
      room for a gain of 10 points (\RoomBar{} or less at body width 4), and \CTwoGained{} of them made it.
      For columns the same count is \COneRoom{} of \COneRoom.
\item \lead{The gain shrank as the creatures grew.} Of the settings with room, all \CTwoRoomFew{} with 16 or
      64 creatures gained 10 points or more, \CTwoRoomMidGained{} of \CTwoRoomMid{} with 256 creatures and
      \CTwoRoomManyGained{} of \CTwoRoomMany{} with 1,024. This fits the sizes: a \BodyTimes-fold wider body
      makes the whole model up to \BodyFewHigh{} times as large with 16 creatures, but at most \BodyManyHigh{}
      times with 1,024, where the word table outweighs everything else.
\item \lead{Twice the columns never did less than a \BodyTimes-fold wider body.} Set the narrowest body (4)
      against the widest (128) at half the columns: the model with twice the columns was ahead in all
      \DoubledMore{} comparisons where the other was not yet at \LearnedBar, and within half a point of it in
      the other \DoubledTied.
\item \lead{The count is not firm.} \CTwoByChanceCap{} of the \CTwoGained{} is a chance case, and with the worst
      repeat at body width 128 set against the best at body width 4 the count is \CTwoRunLow; best against
      worst, \CTwoRunHigh{} (finding~\ref{f:rate}).
\end{itemize}

\finding{Where two models were nearly the same size, the one with more columns knew more}\label{f:size}

More columns add far more trained numbers than a wider body: with everything else the same, 64 columns have
\SizeGridLow{} to \SizeGridHigh{} times the trained numbers of 8 columns, and body width 128 has \BodyGridLow{}
to \BodyGridHigh{} times those of body width 4. The claim says as much: ``there are more weights to help
distinguish the tokens''. So is it the columns, or only the number of trained numbers? With few creatures the
runs can tell: there a model with 8 columns and the widest body is about as large as one with 16 columns and a
narrow body (Figure~\ref{fig:size}).

\fighere{fig09_size}{fig:size}{\point{At nearly the same size, 16 columns with the narrowest body knew more
than 8 columns with the widest; with 1,024 creatures no two models with different columns are near in size.}
A line is one number of table columns; along it the body grows from 4 to 128 neurons, which makes the model
larger. In the middle panel the ringed models differ by \SizeSecondPercent{} percent in size and by
\SizeSecondGap{} points.}

There are \SameSizePairs{} such pairs, in \SameSizeLanguages{} languages: the same bits, creatures and word
use, the model with more columns at most \SameSizePercent{} percent larger (in \SameSizeSmaller{} pairs it is
the smaller), and at least one of the two below \LearnedBar. In all \SameSizeHigher{} the model with 16 columns
scored higher: by 2 points or more in \SameSizeByTwo, by 5 or more in \SameSizeByFive, by as much as
\SameSizeLargest. In \SameSizeAllRuns{} of the \SameSizePairs, every repeat of the model with 16 columns ended
above every repeat of the other.

\lead{The limits are plain.} All \SameSizePairs{} pairs set 8 columns against 16, with 16 to
\SameSizeMostCreatures{} creatures. With 1,024 creatures the word table makes up most of a model, and a model
with twice the columns is always at least \SizeManyClosest{} times as large. And the pairs do not say which of
the parts that grow with the columns did it (section~\ref{s:shows}).

\finding{Up to one bit for each column the facts were known; beyond it, how much a model knew depended on the
creatures and the word use}\label{f:line}

The agent's predictions turn on one number: bits per creature divided by columns, the bits that each number of
a word's row would have to hold. Figure~\ref{fig:loadmap} marks it on the plain dials: in each panel a stepped
line stands where that number is 1. Left of it, a creature has no more bits than its row has numbers.

\fighere{fig10_load_map}{fig:loadmap}{\point{Left of the stepped line the facts are known; right of it they
fall away with many creatures, more so with unequal use.} \term{The typical model}, here and below, is the
middle of the \RunsPerUsePoint{} runs with one word use and one number of columns, bits and creatures. The two
word uses are never put into one middle: where there is much to know their runs lie apart, and a middle of
both would be a score that no run reached. With 16 creatures only the row of 8 columns falls, to
\FewCreaturesLowest{} with unequal use. The \StuckTile{} left of the line is the subject of
finding~\ref{f:plateau}. The runs themselves are on \everyrun.}

\begin{itemize}
\item \lead{Up to the stepped line the facts were known:} the score was \LearnedBar{} or more in
      \AOneFactsKnown{} of the \AOneCases{} settings there. The other \AOneFactsShort{} all have 8 columns, 8
      bits and 1,024 creatures: \AOneShortByKind{} are made by runs one kind of fact short
      (finding~\ref{f:plateau}), and \AOneShortOthers{} missed \LearnedBar{} by \AOneShortOthersBy{} points.
\item \lead{More creatures lowered the score only past the line.} Up to it, 1,024 creatures were known as
      completely as 16. Past it, each fourfold rise in creatures cost points (8 columns, 32 bits, equal use:
      \CreaturesExample). Why, when every creature has a row of its own, these runs do not say.
\item \lead{With 1,024 creatures, each doubling of the columns bought one more doubling of the bits:} the facts
      were known up to as many bits as columns, with both word uses. One doubling further the word uses part:
      \SameLoadUnequal{} with unequal use (8, 16, 32 and 64 columns), but \SameLoadEqual{} with equal use,
      where 64 columns still knew every fact at 128 bits. These heights are not final: \SameLoadRunsLimit{} of
      the \SameLoadRuns{} runs behind them were still training at the limit. With fewer creatures the fall past
      the line is gentler.
\item \lead{A2 held,} within \Limit{} steps. With 1,024 creatures and 4 bits or more per column, none of the
      \ATwoCases{} settings reached \ATwoBar{} (the highest \ATwoHighest; the highest single run
      \ATwoHighestRun). But \ATwoRunsLimit{} of its \ATwoRuns{} runs were ended by the limit, the three of the
      highest setting still gaining \ATwoTopGainLow{} to \ATwoTopGainHigh{} points a doubling.
\item \lead{A1 did not hold as worded; its substance did.} It asked for ``learned'', which needs the questions
      about one creature at \LearnedBar{} too. Of the \AOneCases{} settings at 1 bit per column or less,
      \AOneLearned{} ended as learned (\AOneAllThree{} in all three repeats; \AOneRunsLearned{} of the
      \AOneRuns{} runs). All \AOneMissed{} misses are among the \AOneAtOne{} settings at exactly 1 bit per
      column, \AOneMissedEight{} of them with 8 columns. In \AOneQuestionsShort{} of them
      the facts score was \LearnedBar{} or more, and in \AOneQuestionsOnly{} the questions were the only
      shortfall: the rule had ended those runs because their facts could rise no further
      (finding~\ref{f:rule}).
\item \lead{A3 did not hold.} Of \AThreeCases{} pairs of settings with the same bits per column and otherwise
      alike, \AThreeWithin{} ended within 10 points and \AThreeApart{} did not. In \AThreeMoreHigher{} of the
      \AThreeApart{} the setting with more columns scored higher: the same bits per column are not the same
      load. (\AThreeByChanceCap{} of the \AThreeApart{} are chance cases.)
\end{itemize}

\finding{Some runs stopped exactly one kind of fact short}\label{f:plateau}

A number of runs ended at almost exactly 81.2. With four kinds of fact, \PlateauScore{} is what a model scores
that knows three kinds and guesses the fourth (Figure~\ref{fig:plateau}).

\fighere{fig11_plateau}{fig:plateau}{\point{A run at 81.2 knows three kinds of fact completely and guesses
the fourth.} The \PlateauGroup{} runs of one corner of the grid; a run's score is the average of its four
tiles. With body width 4 the three repeats ended at \PlateauExample.}

\PlateauHereCap{} of these \PlateauGroup{} runs ended one whole kind short, not always the same kind. In the
whole experiment \PlateauRunsWord{} runs did: these \PlateauHere, \PlateauFewer{} more with 256 creatures, and
\PlateauMoreBits{} at 16 bits, where seven kinds known of eight gives \PlateauScoreEight. Such a run stands
still for a long time, so the rule ends it as ``no more gain''. This report calls a setting whose score is made
by such runs a \term{chance case}.

Which repeat it strikes looks like chance: at \PlateauBeside{} of the \PlateauBodies{} body widths in the
figure, a run one kind short stands beside a repeat of the same setting that knows every fact. That it happens
at all is not chance: all \PlateauRunsWord{} have 8 columns and equal use, and none of the \PlateauWiderRuns{}
runs with more columns ended so. \PlateauCarriedCap{} of the \PlateauRunsWord{} were carried on
(finding~\ref{f:rule}): \PlateauCarriedLearned{} then learned the missing kind, \PlateauCarriedNot{} had not
after two more doublings. Why a model holds back a whole kind, these runs do not say.

\finding{Carried on after the rule had ended them, \CarriedGainedDigits{} of \Carried{} runs gained a point or
more, all among the \CarriedRoom{} that stood below \LearnedBar}\label{f:rule}

The plan promised two proofs that the rule does not end runs too early. \lead{First, what the ended runs
were doing.} \NoGainRuns{} runs ended as ``no more gain''. \NoGainKnew{} of them knew the facts and could not
rise. For the other \NoGainBelow, of which \NoGainEight{} have 8 columns, the trainer recorded the gain over
the last doubling: \NoGainMiddle{} points in the middle, more than half a point for \NoGainOverHalf{} of them
and under 1 point for all, as the rule demands. Their training loss (how wrong the model's guesses on its text
still are) fell by \NoGainLossFall{} from about \NoGainLoss. These runs were creeping, not standing still:
nearly level, but not finished.

\fighere{fig12_carried_on}{fig:carried}{\point{Of \Carried{} runs trained on, \CarriedKnew{} had no room to
gain; of the \CarriedRoom{} that had, \CarriedGained{} gained a point or more, and one of the three runs one
kind short learned the missing kind.} That run reached \JumpTo{} some \JumpStepsAfter{} steps after the rule's
stop; it already scored \JumpQuestions{} on the questions about one creature.}

\lead{Second, what the carried-on runs did.} Of the \Marked{} runs marked beforehand, \MarkedNoGain{} ended
as ``no more gain'': \MarkedAtLastLook{} at the last look, where no training was left to add, and \Carried{}
that were then trained on (Figure~\ref{fig:carried}) for two more doublings, \CarriedOneDoubling{} of them for
one, up to the limit. \CarriedKnew{} of the \Carried{} already knew the facts and could not gain a point. Of
the \CarriedRoom{} below \LearnedBar, \CarriedGained{} gained a point or more: \CarriedOthers{} by
\CarriedOthersLow{} to \CarriedOthersHigh{} points and one by \JumpGain, from \JumpFrom{} to \JumpTo. So the
measured error of the rule, as the plan defined it, is \CarriedGainedDigits{} in \Carried, about
\CarriedPercent{} percent; among the runs that could gain, \CarriedGainedDigits{} in \CarriedRoom, only one of
them by much.

\lead{Two things this check does not cover.} First, the rule watches the facts only, and the questions were
often still rising: \CarriedQuestionsGained{} of the \Carried{} runs gained a point or more on questions about
one creature, and of the \CarriedLight{} at 1 bit per column or less, \CarriedLightLearned{} met the bar for
``learned'' after the extra training. So A1's count is lower than longer training would give. Second,
\LimitRuns{} runs were ended by the limit, \LimitStillGaining{} of them still gaining a point or more a
doubling: their scores are what \Limit{} steps gave, not what the models can hold.

\finding{The harder questions gave way before the facts}\label{f:scores}

As in the second experiment, the three kinds of question give way before the plain facts
(Figure~\ref{fig:scores}). With 1,024 creatures, unequal use and as many bits as columns, the typical model
scored \AtLineFactsLow{} on the facts and \AtLineQuestionsLow{} to \AtLineQuestionsHigh{} on questions about
one creature, but \AtLineNewLow{} to \AtLineNewHigh{} on new combinations and \AtLineNotLow{} to
\AtLineNotHigh{} on ``not''. With 64 columns, 128 bits and equal use it scored \HeavyFactsEqual{} on the
facts, but \HeavyQuestionsEqual, \HeavyNewEqual{} and \HeavyNotEqual{} on these.

\fighere{fig13_scores}{fig:scores}{\point{Every part of the test holds for more bits with more columns, and
the three kinds of question fall sooner and further than the facts.} Equal use gives the same order.
Single runs are shown in this report for facts known only.}

\finding{A third of the learning rate changed little; C2's count stayed at its line}\label{f:rate}

The plan provided one more test, at the owner's wish: ``if changing the learning rate leads to better results,
it has to be tested''. After the counted runs, the best repeat of every setting was trained again from the same
starting numbers at a third of the chosen rate, \RcRate, with the same rule for ending: \RcRuns{} runs
(Figure~\ref{fig:ratecheck}). The best repeat is the one that knew most facts; where two tied, as in \RcTies{}
settings, mostly at 100, the questions about one creature decided. Because it is the best of three, the run is
compared in two ways: with its setting's usual score (the middle of three), to see whether the rate helps, and
with itself, to apply the predictions. The plan also allowed a stage at a tenth of the rate, if time on the
rented graphics cards was left. \RcTenth

\fighere{fig14_rate_check}{fig:ratecheck}{\point{At a third of the learning rate nearly every setting ended
within 5 points of its score, and only C2's count changed sides.} Right: a tick where a prediction holds, a
cross where it does not. These counts are over single runs (the best repeat of every setting), so they differ
from those of Figure~\ref{fig:predictions}.}

\begin{itemize}
\item \lead{Better results? Not overall.} \RcWithin{} of the \RcRuns{} runs ended within 5 points of their
      setting's score; \RcHigher{} ended 5 or more higher and \RcLower{} lower. On questions about one creature
      \RcQuestionsHigher{} ended 5 or more higher and \RcQuestionsLower{} lower.
\item \lead{A little worse with few columns, after the same \Limit{} steps.} With 8 columns \RcDownEight{}
      of the \RcPerColumns{} settings ended a point or more below their score and \RcUpEight{} above it; with
      16 columns \RcDownSixteen{} and \RcUpSixteen; with 32 columns \RcDownThirtyTwo{} and \RcUpThirtyTwo;
      with 64 columns \RcDownSixtyFourWord{} and \RcUpSixtyFour, \RcTopSixtyFour{} being at 100 already. Those
      that ended lower lost \RcLostMiddle{} points in the middle. But the lower rate learns more slowly:
      \RcLimitThird{} of its runs were still training at the limit, against \RcLimitChosen, and gaining
      \RcLimitGainThird{} points a doubling against \RcLimitGainChosen. So this is what \Limit{} steps gave,
      not what the lower rate can reach.
\item \lead{The gaps between the column counts did not close.} With 1,024 creatures at 128 bits the runs at a
      third of the rate knew \RcTypicalEqual{} facts with equal use (8, 16, 32 and 64 columns) and
      \RcTypicalUnequal{} with unequal use; the counted runs gave \TypicalEqual, and \TypicalUnequal.
\item \lead{Four verdicts stay; C2's count crosses its line, by less than single runs differ.} C1 and A2
      hold at both rates, A1 and A3 at neither. For C2 the line is \CTwoAllowed. The count is
      \CTwoGained{} by the plan's rule (middles of three), \RcCTwoChosen{} on the best repeats alone, and
      \RcCTwoThird{} on the same repeats at a third of the rate, \RcCTwoByChance{} of them again a chance
      case. But at the chosen rate alone, single runs already give \CTwoSameRepeatLow{} to
      \CTwoSameRepeatHigh{} with one repeat used throughout, and \CTwoSingleLow{} to \CTwoSingleHigh{} with
      any two. So \RcCTwoChosen{} against \RcCTwoThird{} shows no effect of the rate; it shows that C2 sits on
      its line.
\end{itemize}

% ------------------------------------------------------------------------------------------------
\section{What this shows, and what it does not}\label{s:shows}
% ------------------------------------------------------------------------------------------------

\lead{It shows} that with everything else held fixed, more table columns meant more facts known wherever there
was room to improve: what the owner's claim says about the dimensions of a token. A wider body helped much
less, mostly with 8 or 16 columns, and less as the creatures grew. Where sizes could be matched, more columns
beat the same number of trained numbers.

\lead{It does not show:}
\begin{itemize}
\item \lead{Which of the parts that grow with the columns helped:} the rows, the attention parts, the input to
      the middle parts or the output part. If the claim means the row alone, these must be told apart. If it
      means the space the layers work in, they are one thing, and it remains to tell that space from the plain
      number of trained numbers. The owner's words when he rejected a design that kept the space fixed
      (``\dots which maintains the space a model sees'') point to the second; the choice is his.
\item \lead{Whether more columns beat plain size beyond 16 columns, or with 1,024 creatures.} The
      \SameSizePairs{} pairs of finding~\ref{f:size} all set 8 columns against 16, with at most
      \SameSizeMostCreatures{} creatures.
\item \lead{That the low scores are the most these models can hold.} \LimitRuns{} runs were still training
      when \Limit{} steps ended them, gaining \LimitGainLow{} to \LimitGainHigh{} points a doubling in the
      middle. With 1,024 creatures at 128 bits neighbouring column counts lie \TypicalGapLow{} to
      \TypicalGapHigh{} points apart: their order is firm, the exact heights are not.
\item \lead{What a learning rate fitted to each model size would give.} In the practice runs 8 columns did
      best at 0.01 (\PracticeEightMid{} against \PracticeEightLow{} at 0.003) and 64 columns at 0.003
      (\PracticeSixtyFourLow{} against \PracticeSixtyFourMid). Like finding~\ref{f:rate}, this puts the rate's
      effect at a few points, far less than the distance between column counts.
\end{itemize}

\lead{How firm the verdicts on the claim are.} C1 is firm: no fall in \COneComparisons{} comparisons, \COneGained{}
cases where \COneNeeded{} were needed (at least \COneRunLow{} however the repeats are taken), and the same at a
third of the learning rate. C2 is not firm either way: by the rule fixed beforehand it did not hold
(\CTwoGained{} settings, \CTwoAllowed{} allowed), but single runs give \CTwoRunLow{} to \CTwoRunHigh. That a
wider body helps with 8 and 16 columns where there is room is firm; whether in under a tenth of the settings
is not decided.

% ------------------------------------------------------------------------------------------------
\section{What comes next}
% ------------------------------------------------------------------------------------------------

\begin{itemize}
\item \lead{Match sizes where most of the results lie: the agent would do this first.} With 1,024 creatures,
      8 and 16 columns with a body of about \NextBodyLow{} to \NextBodyHigh{} neurons have as many trained
      numbers as a model with twice the columns and the narrowest body. That would say whether more columns
      still beat plain size there. It needs no change to the model, only the owner's word on renting the graphics
      cards again.
\item \lead{Tell apart the parts that grow with the columns.} The owner rejected one way beforehand: a shared
      part that turns a row of any length into one fixed length before the layers read it. Another, for him to
      accept or reject: leave the rows and what the layers read at full length, but let the attention parts
      work at one fixed inner size, as the middle parts already do with the body width.
\item \lead{A rule for ending that also watches the questions.} A1's misses came from runs ended with the facts
      at 100 and the questions still rising, and the rule cannot tell a run one kind short from a finished one.
\item \lead{Runs can be carried on.} Every run kept its saved model, so the \LimitRuns{} runs at the limit and
      the other \PlateauOthers{} runs one kind short can be trained further. A saved model holds the trained
      numbers, not the trainer's running state: the trainer restarts from them.
\end{itemize}

% ------------------------------------------------------------------------------------------------
% Every run: four pages on their side. The heading and its paragraph stand at the foot of the
% upright page before them, so that each wide page is all picture.
% ------------------------------------------------------------------------------------------------
\Needspace{0.85in}
\section{Every run}

The next four pages show all \RunsAll{} runs, a page for each number of table columns: the three repeats of
each of its \SettingsPerColumns{} settings, as facts known when the run ended.

\pageshape{11in}{8.5in}{9.4in}{7.22in}{0.12in}
\fighere{fig15_every_run_8}{fig:everyeight}{\label{page:everyfirst}\point{With 8 columns the facts are known at
8 bits and fall away from 16 bits on, sooner with more creatures and with unequal use.} To find a run, take
the half of its word use, the block of its creatures, the row of its body width and the group of its bits. The
tiles at 81 in the 8-bit groups are runs one kind of fact short (finding~\ref{f:plateau}). A tile at 100 with a
corner mark knew every fact but was still short on the questions about one creature at \Limit{} steps.}

\clearpage
\fighere{fig16_every_run_16}{fig:everysixteen}{\point{With 16 columns the facts are known up to 16 bits, and
with 16 creatures at every number of bits; beyond that the larger vocabularies fall away.} Read as
Figure~\ref{fig:everyeight}.}

\clearpage
\fighere{fig17_every_run_32}{fig:everythirtytwo}{\point{With 32 columns the facts are known up to 32 bits, and
with 16 or 64 creatures at every number of bits.} Read as Figure~\ref{fig:everyeight}. This page contains the
grid of the second experiment, and the two agree: up to 32 bits no setting differs by more than \EarlierApart{}
points (\EarlierSettings{} settings); at 64 bits the third experiment, which could train four times as long,
scores \EarlierHeavyLow{} to \EarlierHeavyHigh{} points higher (\EarlierHeavy{} settings).}

\clearpage
\fighere{fig18_every_run_64}{fig:everysixtyfour}{\label{page:everylast}\point{With 64 columns nearly everything
is known; only 1,024 creatures at 128 bits with unequal use fall clearly short.} Read as
Figure~\ref{fig:everyeight}.}

\pageshape{8.5in}{11in}{6.4in}{9.46in}{-0.12in}

% ------------------------------------------------------------------------------------------------
\appendix
\smallcaps{APPENDIX}\par\nopagebreak\vspace{-9pt}
\section{The learning rate, and how the report counts}
% ------------------------------------------------------------------------------------------------

The one learning rate was fixed from \PracticeRuns{} practice runs: every number of columns and of bits, body
widths 4 and 128, 16 and 1,024 creatures, the rates \PracticeRates, equal use, \PracticeSteps{} steps. The
rule: the rate with the most facts known on average; rates within half a point of the best tie, and the
questions about one creature decide. The two lower rates tied; the questions chose \Rate.

\begin{center}\sffamily\footnotesize
\begin{tabular*}{\linewidth}{@{\extracolsep{\fill}}lrrrrrr@{}}
\toprule
 & \multicolumn{2}{r}{Average over all \PracticeSettings{} practice settings} & \multicolumn{4}{r@{}}{Facts known, by table columns} \\
\cmidrule(lr){2-3}\cmidrule(l){4-7}
Learning rate & facts known & questions about one creature & 8 & 16 & 32 & 64 \\
\midrule
\PracticeRows
\bottomrule
\end{tabular*}
\end{center}

\begin{itemize}
\item \lead{The totals.} C1: \COneCases{} cases (6 body widths × 5 numbers of bits × 4 of
      creatures × 2 word uses), 3 comparisons each. C2: \CTwoCases{} settings (4 column counts ×
      5 × 4 × 2). A1: 10 combinations of columns and bits × 6 × 4 × 2 =
      \AOneCases{} settings. A2: 6 combinations × 6 body widths × 2 word uses = \ATwoCases. A3:
      20 pairs of combinations × 6 × 4 × 2 = \AThreeCases{} pairs.
\item \lead{Word use in C1 and C2.} The plan names body width, bits and creatures as held fixed; this report
      holds the word use fixed too. With the two word uses taken together (the middle of six runs) the
      verdicts are the same: C1 with \PooledCOneGained{} of \PooledCOneCases{} cases, C2 with
      \PooledCTwoGained{} of \PooledCTwoCases{} settings.
\end{itemize}

\section{Where it ran, what it cost, and what differed from the plan}

\begin{itemize}
\item \lead{Where, when and at what cost.} On Colab, a paid cloud service, on 23~November from \StageStart{} to
      \StageEnd{} Pacific time: \Machines{} machines with graphics cards of \MachineTypes{} types
      (\MachineList), up to \MachinesAtOnce{} at once and \MachinesPerType{} of a type, for about \ComputeUnits{}
      of Colab's compute units. The
      counted runs made \StepsMillions{} million training steps; the plan had estimated \PlannedSteps{} million.
\item \lead{Different from the plan:}
  \begin{itemize}
  \item The plan was written for Kaggle's free graphics cards; the owner moved it to Colab for speed.
  \item A group of \RunsPerGroup{} runs trains in one program, and groups were moved between machines: at
        least \GroupsOnSeveralMachines{} of the \Groups{} groups trained on more than one machine and
        \GroupsOnSeveralTypes{} on more than one type of card. At the end \MovedLateGroups{} were moved in
        mid-round to faster machines, saving their state first.
  \item From \SharingStart, several programs shared one card on the G4, A100 and H100 machines (NVIDIA's
        MPS). This changes speed, not arithmetic; \MachinesPaused{} machines were paused to measure it.
  \item The launcher, the owner's program that hands out the groups, was restarted by hand
        \LauncherRestarts{} times to improve it; the machines trained on. \GroupFailuresCap{} group failed
        with a card error and was restarted.
  \item Runs on different types of card, and runs that were moved, are not repeatable bit for bit: another
        card changes the last digits of the arithmetic, which acts like another random start. The repeats show
        how much that matters: in half of all settings the three agree within \SpreadMiddle{} points; in
        \SpreadOverFive{} of the \Settings{} they end more than 5 points apart. The card itself was not measured
        (a setting's repeats trained together), but every column count was trained on all \MachineTypes{} types.
  \item For every run ended as ``no more gain'' the plan asked for its scores at a quarter, a half and the
        whole of its length and the fall of its loss. Finding~\ref{f:rule} gives these in words only.
  \end{itemize}
\item \lead{The rate check} ran from \RcStart{} to \RcEnd{} on \RcMachines{} machines (\RcMachineList); its
      cost is not in the \ComputeUnits{} units. In its first quarter of an hour 34 groups failed, up to three times each, for lack
      of card memory (too many training programs on one card); all were restarted and finished.
\item \lead{As planned:} a run's result is its scores where the rule ended it, for a carried-on run too. Every
      counted run is reported; nothing was rerun to get a better number.
\end{itemize}

\end{document}
